% ID: SEL-D0045-OB-E04
% ESTADO: cotejada. Segunda pasada visual y borrador compilado revisados.
% Fuente: Selectividad/sel_2014_matematicas_aplicadas/EXAMEN -ANDAL-5.pdf
% SHA-256: fe436fc3af8e56d0694316005aedeca2fffefbeb445aa9f60efa0aa6da4cba25
% Página: 2; opción B; ejercicio original 4.

Para estimar la proporción de habitantes que es favorable a la construcción de un centro comercial en un municipio, se ha obtenido el intervalo de confianza $(0.31,0.39)$, al 94\%.
\begin{enumerate}[label=\alph*)]
\item \textbf{(1 punto)} \textquestiondown{}Cuál ha sido el valor de la proporción muestral?
\item \textbf{(0.5 puntos)} Si la muestra aleatoria elegida de esa población para el estudio fue de 500 personas, \textquestiondown{}cuántas de ellas deseaban la construcción del centro comercial?
\item \textbf{(1 punto)} Se desea repetir el estudio para obtener un intervalo de confianza con un error máximo de 0.03 y el mismo nivel de confianza. \textquestiondown{}Cuántas personas, como mínimo, debe tener la nueva muestra aleatoria?
\end{enumerate}
